\section{Introduction}

Long-horizon coding tasks ask an agent to enter an unfamiliar repository, find the cause of a failure, edit, and verify, often across dozens of tool calls.
The default response to failure is capacity: enlarge the window, retrieve more files, compact earlier turns less aggressively, let the agent search for longer.
Each of these assumes that the agent fails because the right material is missing from its prompt or too far away to keep.
That assumption held only partly in our runs.
More material helped, but not in proportion to how much we added, and several failures occurred with the relevant file already in the window: the agent read the failing test, then edited a different function.
The bottleneck we could measure was not how much context the model held, but whether the evidence it needed---where to look, what broke, and when to stop---reached it intact and in a form it could act on.

Our first observation is that direction beats volume.
We ran a paired comparison on a subset of SWE-bench Verified sized at 16 instances\fact{A1}, using DeepSeek V4.1 Flash and a single flag that turns context assembly on or off; task statements were copied verbatim, and the two arms differed in nothing else.
The assembly arm resolved 10/16\fact{A2}; the plain arm resolved 5/16\fact{A3}.
Two instances fail even with the gold patch here, so we also report the set without them: 10/14\fact{A4} against 5/14\fact{A5}, and 10/14\fact{A6} fail-to-pass tests repaired against 6/14\fact{A6}.
The paired outcomes are one-sided: all 5\fact{A7} discordant instances favored assembly and 0\fact{A7} favored the plain arm.
The gain is not uniform across difficulty: in the easiest stratum the two arms resolved 7/7\fact{A8} and 4/7\fact{A8}, in the middle stratum 2/5\fact{A9} and 1/5\fact{A9}, and in the hardest 1/2\fact{A10} and 0/2\fact{A10}.
Nor did it buy its result with volume: mean tokens per instance were 90,431\fact{A11} with assembly and 96,184\fact{A12} without, and because the per-instance pairs are mixed we read this as ``not more expensive'' rather than as a saving.

A second observation concerns instrumentation.
In an earlier ablation on one long task, the locator card was assembled, counted against its budget, and reported as delivered; it had never reached the model.
The renderer dropped the section, and a search of the recorded request bodies for the card's first line found nothing\fact{E9}.
Once the card was actually rendered, first reads landed in the gold region on 4/4\fact{E6} attempts, but only 1/4\fact{E6} resolved.
Knowing where to look is necessary and not sufficient.
We therefore stopped trusting the assembler's own log: every claim below about what the model saw is checked against the request body.

The third observation concerns stopping.
On the same long task with MiniMax-M3.1-Flash, six process constraints written into the system prompt left the resolve count at 0/3\fact{E5}, and per-round output grew rather than shrank.
Prose in the prompt was read and not followed.
Replacing it with a mechanical turn-level reminder produced a different failure: the agent began writing patches but kept producing new versions without declaring one final, 0/2\fact{E7}.
Only a harness rule that made a third revision impossible---at most two patch versions, then submit---closed the task at 5/5\fact{E8}.
Repairs aimed at capacity behaved the same way: evicting repeated reads still gave 0/3\fact{E2}, enlarging the per-kind share 0/3\fact{E3}, and raising the session token ceiling 0/2\fact{E4}.

The fourth observation concerns cost.
Replaying the recorded request bodies of both arms on the same six tasks with MiniMax-M3.1-Flash, we attributed every token.
Our arm emitted 73,051\fact{D2} output tokens per task against 12,495\fact{D2} for the official mcode CLI, while consuming 21,941\fact{D3} fresh input tokens against 52,996\fact{D3}.
The input side is cheaper for us and the output side is not, which reverses the intuition that a context-heavy harness pays for input.
The two cost accountings disagree, and we report both: if every input token is priced as fresh, our arm looks 3$\times$\fact{H1} more expensive, whereas on the actual bill the ratio is 1.66$\times$\fact{D1} (\$0.0486\fact{D1} against \$0.0292\fact{D1} per task).
Spend also concentrates in failure: successful tasks averaged 78,844\fact{D6} tokens, while two failing tasks contributed 70.9\%\fact{D6} of all output.
A single probe that lowered reasoning effort from high to low cut thinking tokens from 941\fact{D4} to 243\fact{D4}, which suggests that the output spend is a knob rather than a property of the task.

The effect is not peculiar to one model family.
On the same six tasks with MiniMax-M3.1-Flash our arm resolved 3/6\fact{C4}; the official CLI resolved 3/6\fact{C1} under its default cap of 24\fact{C1} steps, and 5/6\fact{C1} when that cap was raised to 60\fact{C1} steps.
The two did not solve the same tasks: ours repaired the pytest instance, the CLI the pylint one\fact{C7}.
What we press is narrower and, we think, more useful: the differences we can measure track what reached the model and when it was allowed to stop.

Three contributions follow.
First, a context compiler: the prompt is built rather than concatenated, with typed sections under per-kind share budgets, evidence cards generated from the failing test and the repository graph, and a ledger that records what was rendered.
Second, a controlled measurement of delivery: 10/16\fact{A2} against 5/16\fact{A3} in a paired run, and a cut-by-cut ablation in which prompt rules moved nothing while a mechanical gate moved the same model to 5/5\fact{E8}.
Third, a cost attribution on replayed requests: output rather than input tokens, failure rather than success, and a ratio we report under both accountings, 3$\times$\fact{H1} when all input is priced as fresh and 1.66$\times$\fact{D1} on the bill.

% FACTS-USED: A1, A2, A3, A4, A5, A6, A7, A8, A9, A10, A11, A12, C1, C4, C7, D1, D2, D3, D4, D6, E2, E3, E4, E5, E6, E7, E8, E9, H1
